\section{Future Directions}
\label{sec:future}

\draftsuggestion{Treat foundation-model construction as the end goal while separating tokenizer-specific evidence from model-dependent evidence. Organize the research agenda around three questions: how tokenizers should be evaluated with and without foundation models, what foundation models require from tokenizers, and how the two systems should be designed together. For each direction, state a testable claim, a controlled comparison, and the conditions under which a simpler tokenizer remains preferable.}

\subsection{Evaluating Tokenizers with and without Foundation Models}

\draftsuggestion{Representative foundations: rate--distortion theory \citep{shannon1948}, VQ-VAE \citep{oord2017vqvae}, and the scaling analysis of Byte Latent Transformer \citep{pagnoni2025blt}. Organize evaluation around compression, intelligence, and their joint scaling. Without a foundation model, measure input coverage, encoding stability, token rate, vocabulary or codebook use, reconstruction under the declared equivalence, and the cost of training, encoding, and decoding. Report lossless rate when exact recovery is required, and report distortion together with structural or scientific validity for lossy tokenizers. With a foundation model, test whether the retained information supports prediction, generation, simulation, and control under matched data, backbone, optimization, parameter count, context budget, and total compute. Connect the two levels through rate--distortion--utility curves across model and data scales. Use raw bytes, pixels, residues, atoms, or observations as common denominators because token counts vary across tokenizers. Include reuse as a model-conditioned test: compatibility asks whether models can consume the same token meanings; transferability asks whether a tokenizer remains useful on new data or tasks; replaceability asks how much adaptation is required to substitute one tokenizer for another. Separate tokenizer-training and retokenization costs from foundation-model training, and do not attribute a full-system gain to tokenization without a controlled comparison.}

% Compression asks how efficiently a token representation preserves a declared target. Intelligence asks whether a foundation model can use the preserved information. Neither criterion determines the other. Low reconstruction error can retain details that do not help learning, while a compact representation can remove rare, causal, or task-relevant information. A strong foundation model can also compensate for a weak tokenizer. Future benchmarks should therefore report compression and model utility separately, then study their relation across model sizes and compute budgets.

\subsection{Designing Tokenizers for Foundation Models}

\draftsuggestion{Representative work: \citet{lu2024unifiedio2} and \citet{pertsch2025fast}. Derive tokenizer requirements from the foundation model that consumes its output. Ask how the learning objective, attention pattern, context length, positional representation, output decoder, and deployment budget constrain unit boundaries, token rate, vocabulary or codebook size, and relation encoding. Prediction requires discriminative information and generalization across valid changes. Generation requires sufficient information for reconstruction, valid composition, and controllable output length. Simulation and action further require temporal identity, causal order, and stable closed-loop decoding. Study whether one tokenizer can support these uses or whether shared primitive units should feed task-specific compositions. Require deterministic encoding when models or checkpoints must share token identifiers. Test whether the same tokenizer can serve different foundation-model sizes and architectures, and state the adaptation required when it cannot. For scientific foundation models, add object identity, symmetry, conservation, measurement conditions, and scale as explicit design constraints rather than inferring them from downstream accuracy.}

% Foundation models place different demands on their tokenizers. The relevant unit therefore depends on the model objective and on how the model allocates context and computation. Co-design begins by making these requirements explicit before selecting an algorithmic, neural, or hybrid tokenizer.

\subsection{Co-designing Tokenizers and Foundation Models}

\draftsuggestion{Representative work: \citet{ryoo2021tokenlearner}, \citet{pagnoni2025blt}, and \citet{lu2024unifiedio2}. Reverse the usual design question and ask how the foundation model should change when the tokenizer returns variable-length, overlapping, hierarchical, graph-structured, continuous, or stochastic units. Study positional and relation encodings that respect token geometry and hierarchy; attention and routing that allocate computation according to token scale, uncertainty, or information content; and training objectives that evaluate predictions in the original data space as well as the token space. A model should retain token boundaries, parent--child relations, and decoding uncertainty when these quantities affect the task. Joint training should control codebook collapse, unstable boundaries, and shifts in token meaning, while preserving a clear encoding and decoding procedure. Compare fixed-token backbones with architectures designed for the same tokenizer under matched total compute. Extend checkpoint reuse through token mappings, embedding adaptation, or partial retraining, and report the performance recovery cost. Treat algorithmic tokenizers as valid co-design choices: co-design concerns coordinated assumptions between modules and does not require an end-to-end neural tokenizer.}

% A tokenizer changes sequence length, relation structure, information loss, and the distribution of computation. Foundation-model architectures should expose these choices rather than treating every token stream as a flat sequence of interchangeable identifiers. This reverse direction turns tokenization from preprocessing into a model design variable.

\section{Conclusion}

\draftsuggestion{End the final manuscript with the strongest claim supported by the completed evidence, followed by one clear limitation and one open research question. Remove claims that depend only on unmatched full-model comparisons or unverified tokenizer details.}

% This survey defines a tokenizer through training, encoding, and decoding, then uses that definition to compare algorithmic, neural, and hybrid systems. Separate reviews of AI modalities and domain objects identify reusable component tokenizers. Complex systems add a further problem because these components must be related across modalities, objects, scales, or time.

% The resulting research questions are concrete. Which units preserve information required by a task? Which relations and scientific equivalences must remain explicit? Does the tokenizer improve utility under matched data, model, and compute? Answering these questions requires direct tests of the tokenizer rather than attribution from full-model scores alone.
